% pi_prose_1126.tex -- PI text of 2026-11-26 for paper/journal/v2. It replaces pi_prose_1125.tex.
%
% HOW TO USE (Claude Code)
%   \bn{...}  a number from the build; the value shown is the one printed in draft 4839dea. Replace with its macro.
%   \ct{...}  a citation, by the number it has in draft 4839dea.
%   \chk{...} a clause to check against the build or a source. If it does not hold, send me the clause; do not patch it
%             with a qualifier.
%   Title, abstract, sections 1 (last two paragraphs), 4, 5 and 6 are installed as written. Section 3 takes the two
%   paragraphs given here; the rest of Methods follows the notes at the end.
%   Figure numbers refer to the NEW figure list (see my message): 2 = accuracy against label set, frozen and
%   fine-tuned; 3 = per-class recall (a, b) and pairs of classes (c); 4 = resonance benchmark; 5 = fine-tuning against
%   training size.

\documentclass[11pt,a4paper]{article}
\usepackage[margin=3.2cm]{geometry}
\usepackage[T1]{fontenc}
\usepackage{lmodern}
\usepackage{xcolor}
\usepackage{microtype}
\newcommand{\bn}[1]{\textcolor{blue!70!black}{#1}}
\newcommand{\ct}[1]{\textcolor{black!55}{[#1]}}
\newcommand{\chk}[1]{\textcolor{red!70!black}{#1}}
\newcommand{\bb}{b\bar{b}}
\newcommand{\cc}{c\bar{c}}
\newcommand{\qq}{q\bar{q}}
\title{\Large Label granularity in supervised pretraining for jet tagging\\[4pt]
\normalsize Prose preview, 26 November 2026}
\author{}
\date{}
\begin{document}
\maketitle
\noindent{\small Blue: a number the build supplies. Grey: a citation. Red: a clause to check.}

%%%%%%%%%%%%%%%%%%%%%%%%%%%%%%%%%%%%%%%%%%%%%%%%%%%%%%%%%%%%%%%%%%%%%%%%%%%%%%%%%%%%%%%%%%%%%%%%%%%
% BEGIN ABSTRACT
\begin{abstract}
\noindent
When a model is pretrained by supervised classification, its classes determine which distinctions
the network must learn, yet how fine they need to be remains an open question.
Jets, the collimated sprays of particles produced in high-energy collisions, offer a clean test,
because each simulated jet carries an exact label that can be coarsened without changing the
training data.
We pretrain the Sophon jet-tagging network on the same simulated jets under four nested label
sets, from 188 classes down to 17, each merging classes that differ only in the type of particle
produced, such as bottom and charm quarks.
As the labels coarsen, the frozen representation retains less of what they no longer distinguish:
a linear classifier trained on it identifies the 188 classes with \bn{43\%} accuracy for the
finest labels but \bn{33\%} for the coarsest.
Groupings into 17 classes that cut across this hierarchy still reach \bn{42\%}, so the loss
depends on which classes are merged and not only on how many.
Fine-tuning narrows the gap to about \bn{2} percentage points without closing it, while two
binary benchmarks whose inputs carry little flavor information show no dependence on the label
set.
Fine labels therefore matter most for tasks that rely on the distinctions they preserve, and when
the pretrained model is used without fine-tuning.
An auxiliary regression of the jet mass adds what no label carries, sharpening the mass
resolution of the frozen representation at little cost in accuracy.
\end{abstract}
% END ABSTRACT

%%%%%%%%%%%%%%%%%%%%%%%%%%%%%%%%%%%%%%%%%%%%%%%%%%%%%%%%%%%%%%%%%%%%%%%%%%%%%%%%%%%%%%%%%%%%%%%%%%%
\section{Introduction (last two paragraphs; the first four stay as in draft 4839dea)}
% BEGIN INTRODUCTION TAIL
We evaluate each pretrained model in the two ways such a model is used.
With the network frozen, a linear classifier trained on its features (a linear probe) measures
what the representation holds, both for the 188 pretraining classes and for pairs of classes that
differ in a single flavor.
With the network fine-tuned, we measure what survives adaptation: on the pretraining classes, on
the ten classes of JetClass \ct{8}, which come from a different simulation, and on two public
binary benchmarks, the identification of jets from top quarks (top tagging) \ct{6} and the
discrimination of quark jets from gluon jets \ct{19--21}.
We also repeat the benchmark with which Sophon was introduced, the tagging of a hypothetical heavy
particle $X$ that decays to $\bb$ or to $bs$, with a small classifier trained on the frozen
features.

We find that the frozen representation loses part of what the labels no longer distinguish: the
188 classes are identified less accurately, pairs of classes that a label set merges are separated
less well, and the resonance benchmark loses signal efficiency.
The number of classes does not by itself set this loss, since groupings of the same size that cut
across the hierarchy retain most of the accuracy.
Fine-tuning recovers part of the loss but not all of it, on the pretraining classes and on
JetClass, while the two binary benchmarks, whose inputs carry little flavor information, show no
dependence on the label set.
The mass regression, finally, improves the mass resolution of the frozen features at a small cost
in accuracy.
% END INTRODUCTION TAIL  (no contributions sentence, no numbers here)

%%%%%%%%%%%%%%%%%%%%%%%%%%%%%%%%%%%%%%%%%%%%%%%%%%%%%%%%%%%%%%%%%%%%%%%%%%%%%%%%%%%%%%%%%%%%%%%%%%%
\section{Methods (two paragraphs given; see the notes at the end for the rest)}
% BEGIN METHODS 3.1
A jet is the detector signature of an energetic quark or gluon, or of a heavy particle whose decay
products are collimated by its momentum.
For a machine-learning model it is an unordered set of a few tens of reconstructed particles, each
described here by 17 numbers: its momentum relative to the jet, its electric charge and type and,
for charged particles, the displacement of its track from the collision point, which reveals the
decays of the long-lived hadrons that contain bottom or charm quarks.
Quarks come in several types, or flavors: bottom ($b$), charm ($c$), strange ($s$), and the light
up and down quarks, written together as $q$.
The classes of JetClass-II name what a jet came from.
A resonant class is the decay of a heavy particle into a given set of quarks, gluons and leptons,
each of which appears as one collimated prong of the jet; a QCD class is the ordinary production
of quarks and gluons by the strong interaction, the background in most analyses.
The two binary benchmarks are standard in the field: top tagging asks whether a jet contains the
decay of a top quark, and quark--gluon discrimination whether it was initiated by a quark or by a
gluon.
% END METHODS 3.1

% BEGIN METHODS UNCERTAINTIES (last paragraph of the evaluation subsection; there is no appendix on uncertainties)
Uncertainties are the standard deviation over independently pretrained models, each followed by
one probe or one fine-tuning: three models for each nested label set, two for each model with
mass regression, and one for each of five random groupings.
For quantities that rest on few test jets, such as the background rejection, we also quote the
statistical uncertainty of the test sample, estimated by bootstrap.
% END METHODS UNCERTAINTIES

%%%%%%%%%%%%%%%%%%%%%%%%%%%%%%%%%%%%%%%%%%%%%%%%%%%%%%%%%%%%%%%%%%%%%%%%%%%%%%%%%%%%%%%%%%%%%%%%%%%
% BEGIN RESULTS
\section{Results}

\subsection{The pretraining classes}

Coarser labels leave the frozen features less able to tell the pretraining classes apart
(table~2 and figure~2).
A linear probe on the features of the 188-class model identifies the 188 classes with an accuracy
of \bn{43.5\%}, where chance is \bn{0.5\%} and an untrained network gives \bn{10.2\%}; the
checkpoint released with Sophon gives \bn{43.4\%}, so our retraining reproduces it.
Merging the QCD classes changes little: the 162-class model reaches \bn{43.3\%}, and what it loses
is confined to the QCD classes themselves (table~D4).
The accuracy then falls to \bn{36.1\%} for the 43-class model and to \bn{33.4\%} for the 17-class
model, the two sets that merge flavors within the resonance decays, a loss of \bn{10.1}
percentage points in all.

The loss is not spread evenly over the classes (figure~3(a,b)).
\chk{Classes that a label set leaves untouched are recognized as well as before, while those it
merges with a class of another flavor lose recall, most of all when both their quark and their
lepton flavors are merged.}

The number of classes does not by itself determine the loss.
The random groupings have as many classes as the 17-class set, of the same sizes, yet a linear
probe on their features reaches \bn{42.2\%}, above the 43-class model and only \bn{1.3} points
below the 188-class model.
Their classes cut across the hierarchy, so that nearly every distinction the nested sets remove
still separates some of their classes.
In the other direction the fine labels cost nothing:
\chk{a probe for the 17 classes is as accurate on the features of the 188-class model as on those
of the 17-class model itself (table~D9).}

\subsection{Pairs of classes}

Probes for single pairs of classes show the loss most directly (figure~3(c) and table~D4).
The decays $X\to\bb$ and $X\to\cc$ differ only in the flavor of their quarks and share a class
only in the 17-class set.
A linear probe separates them with an AUC of \bn{0.9977} on the features of the 188-class model
and \bn{0.9902} on those of the 17-class model, a fourfold increase in the error $1-\mathrm{AUC}$,
and the four-quark decays $X\to\bb qq$ and $X\to\cc qq$ behave alike, with \bn{0.9939} against
\bn{0.9716}.
One level up, $X\to bc$ joins $X\to bq$ in a single class from 43 classes on, and its separation
from $X\to bq$ and $X\to cs$ falls from an AUC of \bn{0.941} to \bn{0.932} there and to
\bn{0.885} at 17 classes.
Gluons, merged with the light quarks from 43 classes on, follow the same pattern in the pair
$X\to gg$ against $X\to\qq$.

Two observations complete this picture.
First, merging a pair is not by itself what removes its separation.
The random groupings that happen to place $X\to\bb$ and $X\to\cc$ in one class still separate
them with an AUC of \bn{0.9963}, almost as well as the groupings that keep them apart
(\bn{0.9966}) and far better than the 17-class model, in whose labels no class distinguishes
bottom from charm at all.
Second, a distinction that the labels keep is not fully protected: at 43 classes $X\to\bb$ and
$X\to\cc$ are still separate classes, yet their AUC is already slightly lower than with the
188-class model.

\subsection{Resonance benchmark}

The benchmark with which Sophon was introduced shows what this means for a search (figure~4 and
table~D6).
At a background efficiency of 1\%, a classifier trained on the frozen features of the 188-class
model selects $X\to\bb$ with a signal efficiency of \bn{71.8\%} and $X\to bs$ with \bn{69.2\%}.
With the 17-class model these fall to \bn{65.9\%} and \bn{59.4\%}, and the 43-class model lies in
between, at \bn{70.8\%} and \bn{67.1\%}.
The loss is larger for $X\to bs$, which is a class in none of the label sets and has to be built
entirely from the features.
Even the coarsest labels leave features that are better than none:
\chk{the same classifier without pretraining reaches \bn{55.4\%} for $X\to bs$ [state in Methods
exactly what this reference is]}.
For $X\to\bb$ the two finest label sets also allow Sophon's own selection, a discriminant summed
from the output scores, which reaches \bn{69.8\%} with the 188-class model, close to the trained
classifier.

\subsection{Transfer of the frozen features}

On tasks other than the pretraining classes the frozen features depend much less on the label set
(table~2 and figure~2).
On JetClass a linear probe reaches \bn{77.1\%} with the 188- and 162-class models, \bn{77.0\%}
with the 43-class model and \bn{76.1\%} with the 17-class model.
The loss of one point comes mainly from the decays of the Higgs boson to charm and to bottom
quarks, whose recall falls by \bn{5.0} and \bn{2.9} points (table~D7); these are the JetClass
classes that differ in a flavor the 17-class set has merged.
The random groupings, which kept most of the accuracy on the pretraining classes, do not keep it
here: at \bn{75.4\%} they lie below every nested set.

On the two binary benchmarks the label sets cannot be told apart, and on top tagging the small
differences run the other way: the 17-class model has the highest accuracy of the nested sets,
\bn{91.5\%} against \bn{91.1\%} for the 188-class model.
These frozen results are a linear readout of features computed from incomplete inputs
(section~3) and lie far below the fine-tuned ones.

\subsection{Fine-tuning}

Fine-tuning narrows the gap between the label sets without closing it (table~3 and figure~2).
When the whole network is trained on $10^6$ labeled jets of the pretraining classes, the 17-class
model reaches \bn{41.0\%} and the 188-class model \bn{43.3\%}, so the \bn{10.1} points that
separated their frozen features shrink to \bn{2.3}.
The 43-class model recovers to within \bn{1.0} points and the 162-class model fully.
All of them remain far above the same network trained from scratch on these jets, at
\bn{37.3\%}: of the \bn{6.0} points that pretraining provides, the coarsest labels give up more
than a third.
\chk{[One sentence on the two finest sets ending slightly below their own linear probes,
\bn{43.3\%} against \bn{43.5\%} for the 188-class model, written when the two checks end.]}

The ordering is the same on JetClass, whose classes the pretraining never saw.
The 188-class model reaches \bn{85.0\%} and the 17-class model \bn{84.4\%}, a difference of
\bn{0.5} points where pretraining gains \bn{2.1} over the network trained from scratch.
After fine-tuning the loss is confined to the decay of the Higgs boson to bottom quarks
(table~D7).
The random groupings end level with the 43-class set on the pretraining classes, at \bn{42.3\%},
and between the 43- and 17-class sets on JetClass, at \bn{84.7\%}.

The label set matters more when labeled jets are scarce (figure~5).
On the pretraining classes the 17-class model trails the 188-class model by \bn{7.0} points when
fine-tuned on $10^4$ jets and by \bn{4.0} points on $10^5$, before the difference reaches
\bn{2.3} points at $10^6$.

Top tagging and quark--gluon discrimination show no dependence on the label set.
After fine-tuning on the full training sets the four nested sets agree to within \bn{0.04} points
in top-tagging accuracy and \bn{0.06} points in quark--gluon accuracy, and their background
rejections agree within the uncertainties (table~3); the same holds for quark--gluon jets
simulated with Herwig instead of Pythia (table~D2).
Pretraining itself does help: the fine-tuned 188-class model reaches \bn{94.00\%} on top tagging
against \bn{93.87\%} for the network trained from scratch, and \bn{84.33\%} against \bn{83.97\%}
on quark--gluon discrimination.
The corresponding rise in top-tagging rejection, from \bn{1,412} to \bn{1,799}, is less certain
than the spread between models suggests, because only about a hundred background jets pass this
working point, which adds a statistical uncertainty of about \bn{120}.

\subsection{Jet mass}

The auxiliary regression makes the jet mass much easier to recover from the frozen features
(table~4).
An MLP probe on the features of the 162-class model estimates the mass with a resolution of
\bn{21.7\%}, against \bn{24.5\%} for the reconstructed soft-drop mass itself; with the regression
in pretraining the same probe reaches \bn{15.7\%}, and the regression output of the network
\bn{14.7\%}.
The gain is the same at 17 classes, from \bn{20.8\%} to \bn{15.6\%}.
Without the regression, the coarser labels already give a slightly better mass resolution than
the finer ones, \bn{20.8\%} for the 17-class model against \bn{21.8\%} for the 188-class model.

The price in classification is small.
With frozen features the accuracy on the pretraining classes falls by \bn{0.4} points at 162
classes and by \bn{0.3} points at 17, within the variation between models.
After fine-tuning, the models pretrained with the regression are lower by about half a point on
the pretraining classes and by \bn{0.1} points on JetClass, at both label sets.
On top tagging the regression helps the frozen features: the accuracy of the linear probe rises
from \bn{91.0\%} to \bn{91.5\%} at 162 classes and from \bn{91.5\%} to \bn{92.0\%} at 17
(table~2).
% END RESULTS

%%%%%%%%%%%%%%%%%%%%%%%%%%%%%%%%%%%%%%%%%%%%%%%%%%%%%%%%%%%%%%%%%%%%%%%%%%%%%%%%%%%%%%%%%%%%%%%%%%%
% BEGIN DISCUSSION
\section{Discussion}

Taken together, these results show a representation shaped by the distinctions its labels draw.
When a label set stops separating bottom from charm quarks, or gluons from light quarks, the
frozen features lose part of that separation, and the loss appears wherever the distinction is
needed: in the accuracy on the pretraining classes, in the separation of single pairs, and in the
efficiency of a resonance search.
What matters is that the distinction is required somewhere, not that every pair of classes is kept
apart.
The random groupings merge some pairs of flavors and still separate them, whereas the 17-class
set separates them nowhere.
This suggests that a network learns a distinction when any class boundary demands it, and then
applies it to decays whose own labels do not.

The same pattern has been seen in image classification, where coarser labels cost more for frozen
features than after fine-tuning \ct{3}, and classes merged along a semantic hierarchy remain more
entangled than classes merged at random \ct{23}.
What simulation adds is control: the labels are exact, the training examples and their order are
identical for every label set, and each merge removes a distinction that can be named, so the loss
can be traced to individual pairs of classes instead of being inferred from an average over a
taxonomy.

For the use of a pretrained tagger, the consequence depends on how it is used.
Read out from frozen features, or through class scores summed into a discriminant, the model can
offer only the separations its labels asked for, and the classes an analysis will need to tell
apart should be distinct in the pretraining labels.
Fine-tuning relaxes this requirement without removing it: with a million labeled jets the
coarsest labels still give up more than a third of what pretraining provides on the pretraining
classes and a quarter on JetClass, and more when labeled jets are few.
Tasks that do not draw on the merged distinctions are indifferent to them, as top tagging and
quark--gluon discrimination show; in both, the inputs that carry flavor are largely absent.
A grouping that preserves the pretraining classes is not thereby a good label set, however.
The frozen features of the random groupings transfer less well than those of any nested set,
which may reflect that their classes mix decays of different structure.

The mass regression shows that what the labels lack can be supplied by other means: it improves
the mass resolution of the frozen features by about a quarter and costs at most half a point of
accuracy.
Our study has limits: it rests on one network, with \bn{2.2}~million parameters, and on one
simulated dataset.
Each step of the hierarchy removes several distinctions at once, so the cost of a single one is
established only through pairs of classes.
The random groupings were chosen among random draws to overlap little with the hierarchy, and we
have not studied typical draws.
Our fine-tuning uses one recipe for all models, and on top tagging it stays below the best
published results (appendix~F), so the absolute performance on the benchmarks should not be read
as the reach of the method.
% END DISCUSSION

% BEGIN CONCLUSION
\section{Conclusion}

We have varied the labels of a supervised pretraining while keeping its data fixed, which
simulation makes possible and which, for jets, had not been done.
The representation retains what the labels distinguish.
Merging flavors in the labels removes part of their separation from the frozen features, by ten
percentage points of accuracy on the pretraining classes between the finest and the coarsest set,
while a grouping of the same size that cuts across the hierarchy loses only about one.
Fine-tuning recovers most of the difference on the pretraining classes but not all of it, and
tasks that do not need the merged distinctions do not notice them.
Pretraining labels should therefore keep apart the classes that later analyses will need to
separate, above all when the model is to be used without fine-tuning, and quantities that no
label carries, such as the jet mass, can be added as auxiliary targets.
% END CONCLUSION

%%%%%%%%%%%%%%%%%%%%%%%%%%%%%%%%%%%%%%%%%%%%%%%%%%%%%%%%%%%%%%%%%%%%%%%%%%%%%%%%%%%%%%%%%%%%%%%%%%%
\section*{Notes for the rest of Methods}
\begin{itemize}\itemsep2pt
\item Subsections: Jets and classes; Data and model; Label sets; Random groupings and mass
regression; Probes and fine-tuning; Downstream tasks and metrics. The evaluation text is now one
block of seven paragraphs; split it in two at the paragraph on JetClass.
\item ``109,938,036 training jets'' becomes ``about 110 million''. Exact counts go to the appendix.
\item The random groupings are described once, here, in two sentences: random assignments matched
in size to the 17-class set and chosen, among such draws, to share few merges with it. No other
section repeats how they were chosen, except the one sentence in the Discussion.
\item The resonance paragraph says once what ``trained from scratch'' is on that benchmark, and
that our protocol differs from Sophon's. Results does not repeat either.
\item The sentence on the output layer and its bias refit goes to the appendix.
\end{itemize}

\end{document}
